\chapter*{阅读约定：对象、符号与论证层级}
\addcontentsline{toc}{chapter}{阅读约定：对象、符号与论证层级}
\label{front:reading_conventions}
我们首先区分四类内部变量：当前状态 $\mathbf{x}$、持久结构 $S$、历史及中间记忆 $\mathbf{m}$、目标与价值状态 $\mathbf{v}$。输入为 $\mathbf{u}$，控制为 $\mathbf{c}$，动作为 $\mathbf{a}$。符号的具体维度与空间在各模型中声明，不能仅凭名称执行运算。

\begin{table}[htbp]\centering
\begin{tabularx}{\textwidth}{l X}
\toprule 对象 & 使用约定 \\\midrule
状态强度 & 由范数或指定统计量测量，不自动具有能量单位。\\
信息熵 & 对指定概率分布定义；特征向量本身不是概率。\\
认知温度 & 若用于 softmax，表示尺度参数；若用于相图，须重新定义观测量。\\
稳定性 & 区分有界、平衡点稳定、渐近收敛及任务成功。\\
相位与干涉 & 可作为周期变量和带符号耦合的计算机制，不预设量子实现。\\
几何与拓扑 & 必须声明空间、度量、边界与离散化，不由类比自动获得。\\
\bottomrule\end{tabularx}\end{table}

在使用连续时间模型时，我们研究的是计算状态的连续近似；实际系统可以采用离散步、事件触发或异步更新。各尺度的步长和延迟需要单独给出。

\section*{插图与历史符号的解释}
我们保留原有插图及其可编辑源文件。图中的“波”“场”“势能”“宏观意志”分别可用于辅助理解状态传播、分布式表示、目标偏置与元控制；它们不构成相应物理规律在认知系统中成立的证据。几何示意图不保证空间真的满足图中描绘的正交、距离或曲率。

部分历史图和专门模型仍使用 $\Psi$ 表示认知活动场。该符号只在图或局部模型中有效；一般模型使用 $\mathbf{x}$ 或 $\phi$，AgentOS 的自我使用 $\Psi_{self}$，并可在 AgentOS 专用语境中简写为 $\Psi$。历史的 TCE 控制方程与 AgentOS 的 Temperature \& Control Engine 不使用同一对象身份。

\section*{证明、模型与实现}
我们仅在条件明确且推导成立时赋予命题数学保证。边界条件、噪声、离散步长或结构更新发生变化后，原结论需要重新检查。专门模型中的“守恒”“收敛”和“等价”不能直接推广到任意智能系统。物理载体的功耗、散热和通信延迟作为外部实现约束测量。
